% GENERATED FILE. Edit canonical lessons, then run npm run generate:books.
% canonical-source-hash: fnv1a64:580db4c4c7a1a8ae
% canonical-lessons: PA-C44-duja, PA-C44-tija, PA-C44-chautha, PA-C44-pahila, PA-C44-panjvan, PA-R44-first-to-fifth

\chapter{Which One in the Line}
\label{ch:pa-ordinals}
\hlchaptermodality{44}

\begin{chapteropening}
\textbf{By the end of this chapter:} \emph{I can say first to fifth in Punjabi, agree an ordinal with a masculine or feminine thing, and say which of the five endings is the only one that still works on a new number.}

Complete PA-R44-first-to-fifth and run pahilaa to panjvaan, then name the one ending that still works on a number the book has not taught.
\end{chapteropening}

\section[dūjā]{\pa{ਦੂਜਾ} (dūjā) — second}

\label{lesson:PA-C44-duja}



\begin{quote}
A second kind of number word starts here, so the recalls come first, at the four growing distances this book measures.

\begin{itemize}
\raggedright
  \item \emph{Recall:} ask \emph{kinne?} — \textbf{R1}, one lesson back
  \item \emph{Recall:} ask \emph{kauṇ?} — \textbf{R2}, five lessons back
  \item \emph{Recall:} join two halves with \emph{ki} — \textbf{R3}, twenty lessons back
  \item \emph{Recall:} name the cue that says where the age goes on a form — \textbf{R4}, eighty lessons back
\end{itemize}
\end{quote}

\begin{cousinweb}
\begin{quote}
\textbf{\pa{ਦੂਜਾ}} — \emph{dūjā} — \textbf{second}
\end{quote}

Say \textbf{\pa{ਦੋ}} (\emph{do}), then \textbf{\pa{ਦੂਜਾ}} (\emph{dūjā}). The number has not gone anywhere: the \textbf{\pa{ਦ}} is still at the front and the vowel has only lengthened.
\end{cousinweb}

\subsection*{Script — one tiny step}
Punjabi keeps \textbf{two sets of number words}, and they answer two different questions:

\noindent
\begin{tabularx}{\linewidth}{@{}>{\raggedright\arraybackslash}X>{\raggedright\arraybackslash}X>{\raggedright\arraybackslash}X@{}}
\toprule
\textbf{the question} & \textbf{the set} & \textbf{example} \\
\midrule
\textbf{\pa{ਕਿੰਨੇ}?} — how many? & \textbf{\pa{ਇੱਕ} \pa{ਦੋ} \pa{ਤਿੰਨ} \pa{ਚਾਰ} \pa{ਪੰਜ}} & \textbf{\pa{ਦੋ} \pa{ਰੋਟੀ}} — two rotis \\
\emph{which one in the line?} & \textbf{\pa{ਦੂਜਾ}} and its fellows & \textbf{\pa{ਦੂਜੀ} \pa{ਰੋਟੀ}} — the second roti \\
\bottomrule
\end{tabularx}

\textbf{\pa{ਕਿੰਨੇ}}, taught one lesson ago, is exactly the question the first set answers. This chapter teaches the other set.

An ordinal is an ordinary \textbf{-\pa{ਾ}} adjective, so it agrees with what it counts: \textbf{\pa{ਦੂਜਾ} \pa{ਦੋਸਤ}} (the second friend), \textbf{\pa{ਦੂਜੀ} \pa{ਰੋਟੀ}} (the second roti).

\subsection*{Your turn}
Say these aloud:

\begin{itemize}
\raggedright
  \item \emph{do}, then \emph{dūjā}
  \item \emph{dūjī roṭī} — the second roti
  \item which of the two sets \emph{panj} belongs to — \textbf{the \pa{ਕਿੰਨੇ} set}
\end{itemize}

\begin{tcolorbox}[breakable,colback=teal!4,colframe=teal!35!black,title={Before you move on}]
How many sets of number words does Punjabi keep? (\textbf{Two.}) Which number can you still see inside \textbf{\pa{ਦੂਜਾ}}? (\textbf{\pa{ਦੋ}}.)

Sources: \href{https://en.wiktionary.org/wiki/\%E0\%A8\%A6\%E0\%A9\%82\%E0\%A8\%9C\%E0\%A8\%BE}{Wiktionary: \pa{ਦੂਜਾ}}.
\end{tcolorbox}
\section[tījā]{\pa{ਤੀਜਾ} (tījā) — third}

\label{lesson:PA-C44-tija}



\begin{quote}
Four recalls, at four distances.

\begin{itemize}
\raggedright
  \item \emph{Recall:} say \emph{dūjā} — \textbf{R1}, one lesson back
  \item \emph{Recall:} ask \emph{kadoṁ?} — \textbf{R2}, five lessons back
  \item \emph{Recall:} say \emph{maiṁ sochdā hāṁ ki …} — \textbf{R3}, twenty lessons back
  \item \emph{Recall:} name the cue that says where the age goes on a form — \textbf{R4}, eighty-one lessons back
\end{itemize}
\end{quote}

\begin{cousinweb}
\begin{quote}
\textbf{\pa{ਤੀਜਾ}} — \emph{tījā} — \textbf{third}
\end{quote}

Put the two you now have beside the numbers under them:

\noindent
\begin{tabularx}{\linewidth}{@{}>{\raggedright\arraybackslash}X>{\raggedright\arraybackslash}X@{}}
\toprule
\textbf{how many} & \textbf{which one in the line} \\
\midrule
\textbf{\pa{ਦੋ}} (\emph{do}) & \textbf{\pa{ਦੂਜਾ}} (\emph{dūjā}) \\
\textbf{\pa{ਤਿੰਨ}} (\emph{tinn}) & \textbf{\pa{ਤੀਜਾ}} (\emph{tījā}) \\
\bottomrule
\end{tabularx}

The first letter survives in both, and the same \textbf{-\pa{ਜਾ}} sits on the end of each.
\end{cousinweb}

\begin{grammarlens}[title={What you've built}]
That shared \textbf{-\pa{ਜਾ}} is not a Punjabi invention and it is not a rule you can use. Wiktionary traces \textbf{\pa{ਦੂਜਾ}} to Sanskrit \emph{dvitīya} and \textbf{\pa{ਤੀਜਾ}} to Sanskrit \emph{tṛtīya} — one Sanskrit ending, \textbf{-tīya}, worn down to \textbf{-\pa{ਜਾ}} in two words that travelled together.

Two words is a resemblance, not a pattern, and this book will not call it one. Watch what the next number does with it.
\end{grammarlens}

\subsection*{Your turn}
Say these aloud:

\begin{itemize}
\raggedright
  \item \emph{dūjā}, then \emph{tījā}
  \item the number inside each — \emph{do}, \emph{tinn}
  \item \emph{tījī roṭī} — the third roti
\end{itemize}

\begin{tcolorbox}[breakable,colback=teal!4,colframe=teal!35!black,title={Before you move on}]
What do \textbf{\pa{ਦੂਜਾ}} and \textbf{\pa{ਤੀਜਾ}} share? (\textbf{The ending -\pa{ਜਾ}}.) Is that an ending you can put on any number? (\textbf{Not yet shown — two words is a resemblance.})

Sources: \href{https://en.wiktionary.org/wiki/\%E0\%A8\%A4\%E0\%A9\%80\%E0\%A8\%9C\%E0\%A8\%BE}{Wiktionary: \pa{ਤੀਜਾ}}; \href{https://en.wiktionary.org/wiki/\%E0\%A8\%A6\%E0\%A9\%82\%E0\%A8\%9C\%E0\%A8\%BE}{Wiktionary: \pa{ਦੂਜਾ}}.
\end{tcolorbox}
\section[chauthā]{\pa{ਚੌਥਾ} (chauthā) — fourth}

\label{lesson:PA-C44-chautha}



\begin{quote}
Four recalls, at four distances.

\begin{itemize}
\raggedright
  \item \emph{Recall:} say \emph{tījā} — \textbf{R1}, one lesson back
  \item \emph{Recall:} write \textbf{\pa{ਥ}} — \textbf{R2}, five lessons back
  \item \emph{Recall:} write \textbf{\pa{ਕਿ}} — \textbf{R3}, twenty lessons back
  \item \emph{Recall:} say how much space a written answer needs on a form — \textbf{R4}, eighty lessons back
\end{itemize}
\end{quote}

\begin{cousinweb}
\begin{quote}
\textbf{\pa{ਚੌਥਾ}} — \emph{chauthā} — \textbf{fourth}
\end{quote}

\textbf{\pa{ਚਾਰ}} (\emph{chār}) gives \textbf{\pa{ਚੌ}-} (\emph{chau-}), and then the word ends in \textbf{-\pa{ਥਾ}}.

The \textbf{\pa{ਥ}} in it is the letter you wrote five lessons ago for \textbf{\pa{ਕਿੱਥੇ}}. It is doing the same shape here in a word that has nothing to do with \emph{where}.
\end{cousinweb}

\begin{grammarlens}[title={What you've built}]
Three ordinals now, and two endings:

\noindent
\begin{tabularx}{\linewidth}{@{}>{\raggedright\arraybackslash}X>{\raggedright\arraybackslash}X>{\raggedright\arraybackslash}X@{}}
\toprule
\textbf{how many} & \textbf{which one in the line} & \textbf{ending} \\
\midrule
\textbf{\pa{ਦੋ}} & \textbf{\pa{ਦੂਜਾ}} & \textbf{-\pa{ਜਾ}} \\
\textbf{\pa{ਤਿੰਨ}} & \textbf{\pa{ਤੀਜਾ}} & \textbf{-\pa{ਜਾ}} \\
\textbf{\pa{ਚਾਰ}} & \textbf{\pa{ਚੌਥਾ}} & \textbf{-\pa{ਥਾ}} \\
\bottomrule
\end{tabularx}

So \textbf{-\pa{ਜਾ}} was a resemblance and not a rule, exactly as the last lesson refused to promise. Wiktionary says why: \textbf{\pa{ਦੂਜਾ}} and \textbf{\pa{ਤੀਜਾ}} come from Sanskrit \emph{dvitīya} and \emph{tṛtīya}, and \textbf{\pa{ਚੌਥਾ}} comes from Sanskrit \emph{caturtha} — a different Sanskrit ending, worn down its own way.

These are not built words. They are \textbf{inherited} ones, each worn separately, which is why no rule joins them.
\end{grammarlens}

\subsection*{Your turn}
Say these aloud:

\begin{itemize}
\raggedright
  \item \emph{dūjā tījā chauthā}
  \item which two share an ending — \emph{\textbf{dūjā}} and \emph{\textbf{tījā}}
  \item \emph{chauthā dost} — the fourth friend
\end{itemize}

\begin{tcolorbox}[breakable,colback=teal!4,colframe=teal!35!black,title={Before you move on}]
How many endings do the three ordinals use? (\textbf{Two}.) Are they built from the numbers, or inherited? (\textbf{Inherited} — \textbf{each worn down its own way}.)

Sources: \href{https://en.wiktionary.org/wiki/\%E0\%A8\%9A\%E0\%A9\%8C\%E0\%A8\%A5\%E0\%A8\%BE}{Wiktionary: \pa{ਚੌਥਾ}}.
\end{tcolorbox}
\section[pahilā]{\pa{ਪਹਿਲਾ} (pahilā) — first}

\label{lesson:PA-C44-pahila}



\begin{quote}
Four recalls, at four distances.

\begin{itemize}
\raggedright
  \item \emph{Recall:} say \emph{chauthā} — \textbf{R1}, one lesson back
  \item \emph{Recall:} ask \emph{kitthe?} and write it — \textbf{R2}, five lessons back
  \item \emph{Recall:} ask \emph{kyoṁ?} — \textbf{R3}, twenty lessons back
  \item \emph{Recall:} say how much space a written answer needs on a form — \textbf{R4}, eighty-one lessons back
\end{itemize}
\end{quote}

\begin{cousinweb}
\begin{quote}
\textbf{\pa{ਪਹਿਲਾ}} — \emph{pahilā} — \textbf{first}
\end{quote}

Now go looking for the number. \textbf{\pa{ਦੂਜਾ}} kept its \textbf{\pa{ਦ}}. \textbf{\pa{ਤੀਜਾ}} kept its \textbf{\pa{ਤ}}. \textbf{\pa{ਚੌਥਾ}} kept its \textbf{\pa{ਚ}}. This word keeps nothing of \textbf{\pa{ਇੱਕ}} — not a letter.
\end{cousinweb}

\begin{grammarlens}[title={What you've built}]
Four ordinals, three endings, and one of the four with no number in it at all:

\noindent
\begin{tabularx}{\linewidth}{@{}>{\raggedright\arraybackslash}X>{\raggedright\arraybackslash}X>{\raggedright\arraybackslash}X@{}}
\toprule
\textbf{how many} & \textbf{which one in the line} & \textbf{the number, still visible?} \\
\midrule
\textbf{\pa{ਇੱਕ}} & \textbf{\pa{ਪਹਿਲਾ}} & \textbf{no} \\
\textbf{\pa{ਦੋ}} & \textbf{\pa{ਦੂਜਾ}} & yes \\
\textbf{\pa{ਤਿੰਨ}} & \textbf{\pa{ਤੀਜਾ}} & yes \\
\textbf{\pa{ਚਾਰ}} & \textbf{\pa{ਚੌਥਾ}} & yes \\
\bottomrule
\end{tabularx}

That is why \emph{first} arrives fourth in this chapter rather than opening it. A word can only be seen to BREAK a shape after the shape has been seen holding, and it had to hold three times before this one could be read as a surprise.

Wiktionary explains the surprise: \textbf{\pa{ਪਹਿਲਾ}} goes back through Prakrit to Sanskrit \emph{prathama}, which is \emph{\textbf{pra-}} \textquotedblleft{}forward\textquotedblright{} plus a superlative ending. It never meant \emph{one-th}. It meant \textbf{foremost} — and English \emph{first} is the superlative of \emph{fore}, which is the identical idea reached on a different road.
\end{grammarlens}

\subsection*{Your turn}
Say these aloud:

\begin{itemize}
\raggedright
  \item \emph{pahilā dūjā tījā chauthā} — the four, in order
  \item which one has no number inside it — \emph{\textbf{pahilā}}
  \item \emph{pahilī roṭī} — the first roti
\end{itemize}

\begin{tcolorbox}[breakable,colback=teal!4,colframe=teal!35!black,title={Before you move on}]
Which number is inside \textbf{\pa{ਪਹਿਲਾ}}? (\textbf{None} — \textbf{\pa{ਇੱਕ}} is not in it.) Why did the chapter not begin with it? (\textbf{Nothing had been seen for it to break yet.})

Sources: \href{https://en.wiktionary.org/wiki/\%E0\%A8\%AA\%E0\%A8\%B9\%E0\%A8\%BF\%E0\%A8\%B2\%E0\%A8\%BE}{Wiktionary: \pa{ਪਹਿਲਾ}}.
\end{tcolorbox}
\section[panjvāṁ]{\pa{ਪੰਜਵਾਂ} (panjvāṁ) — fifth}

\label{lesson:PA-C44-panjvan}



\begin{quote}
Four recalls, at four distances.

\begin{itemize}
\raggedright
  \item \emph{Recall:} say \emph{pahilā} — \textbf{R1}, one lesson back
  \item \emph{Recall:} ask \emph{kinne?} — \textbf{R2}, five lessons back
  \item \emph{Recall:} answer a \emph{kyoṁ} with \emph{kyoṁki} — \textbf{R3}, twenty lessons back
  \item \emph{Recall:} say how much space a written answer needs on a form — \textbf{R4}, eighty-two lessons back
\end{itemize}
\end{quote}

\begin{cousinweb}
\begin{quote}
\textbf{\pa{ਪੰਜਵਾਂ}} — \emph{panjvāṁ} — \textbf{fifth}
\end{quote}

Say the four you have, then this one, and listen to the end of each:

\begin{quote}
\emph{pahilā · dūjā · tījā · chauthā} … \emph{panjvāṁ}
\end{quote}

The first four end in a plain \textbf{-\pa{ਾ}}. This one ends in \textbf{-\pa{ਵਾਂ}} — the \textbf{\pa{ਾਂ}} carries a nasal, and you can hear it. The seam in this set is audible before it is explained.
\end{cousinweb}

\begin{grammarlens}[title={What you've built: an inherited word that turned into a rule}]
Wiktionary says two things about this word, and it is worth reading both, because they look like they contradict each other and do not.

\begin{enumerate}
\raggedright
  \item \textbf{\pa{ਪੰਜਵਾਂ}} is \textbf{inherited} — from Prakrit \emph{paṃcama}, from Sanskrit \emph{pañcama}. Like the other four, it came down whole.
  \item \textbf{\textquotedblleft{}By surface analysis\textquotedblright{}}, it is \textbf{\pa{ਪੰਜ}} + \textbf{-\pa{ਵਾਂ}}.
\end{enumerate}

Both are true, and the second is what happened AFTER the first. The word arrived in one piece; then speakers looked at it, saw their own word for five sitting at the front, took the rest to be an ending, and started using that ending on other numbers.

So the seam in this chapter is not where Punjabi stopped inheriting. \textbf{Every one of these five words was inherited.} The seam is where the inheriting stopped being visible — where a whole word came to look like a sum, and the language believed it.

That is why \textbf{-\pa{ਵਾਂ}} goes on numbers this book has not taught you yet, and \textbf{-\pa{ਜਾ}} and \textbf{-\pa{ਥਾ}} go on nothing at all.
\end{grammarlens}

\subsection*{Your turn}
Say these aloud:

\begin{itemize}
\raggedright
  \item \emph{panj}, then \emph{panjvāṁ}
  \item the five in order, and listen for where the ending changes
  \item which of the five endings can be put on a new number — \emph{\textbf{-vāṁ}}
\end{itemize}

\begin{tcolorbox}[breakable,colback=teal!4,colframe=teal!35!black,title={Before you move on}]
What can you hear at the end of \emph{panjvāṁ} that the other four do not have? (\textbf{A nasal}.) Were the four before it built or inherited? (\textbf{Inherited — and so was this one.})

Sources: \href{https://en.wiktionary.org/wiki/\%E0\%A8\%AA\%E0\%A9\%B0\%E0\%A8\%9C\%E0\%A8\%B5\%E0\%A8\%BE\%E0\%A8\%82}{Wiktionary: \pa{ਪੰਜਵਾਂ}}.
\end{tcolorbox}
\section[pahilā · dūjā · tījā · chauthā · panjvāṁ]{First to fifth, and the word that turned into a rule}

\label{lesson:PA-R44-first-to-fifth}



\begin{quote}
Four recalls, at four distances, before the run.

\begin{itemize}
\raggedright
  \item \emph{Recall:} say \emph{panjvāṁ} — \textbf{R1}, one lesson back
  \item \emph{Recall:} say how many sets of number words Punjabi keeps — \textbf{R2}, five lessons back
  \item \emph{Recall:} open a condition with \emph{je} — \textbf{R3}, twenty lessons back
  \item \emph{Recall:} say what to do when a written form goes wrong — \textbf{R4}, eighty lessons back
\end{itemize}
\end{quote}

\subsection*{Your turn}
Say these aloud:

\begin{itemize}
\raggedright
  \item the run — \emph{pahilā dūjā tījā chauthā panjvāṁ}
  \item each number with its ordinal — \emph{ikk pahilā}, \emph{do dūjā}, and on
  \item which one of the five ends in a nasal — \emph{\textbf{panjvāṁ}}
  \item \emph{dūjī roṭī}, then \emph{dūjā dost} — the same ordinal on a feminine and a masculine thing
\end{itemize}

\begin{grammarlens}[title={What you've built}]
Four endings across five words:

\noindent
\begin{tabularx}{\linewidth}{@{}>{\raggedright\arraybackslash}X>{\raggedright\arraybackslash}X>{\raggedright\arraybackslash}X@{}}
\toprule
\textbf{ordinal} & \textbf{ending} & \textbf{comes from} \\
\midrule
\textbf{\pa{ਪਹਿਲਾ}} & \textbf{-\pa{ਲਾ}} & Sanskrit \emph{prathama}, \textquotedblleft{}foremost\textquotedblright{} \\
\textbf{\pa{ਦੂਜਾ}} & \textbf{-\pa{ਜਾ}} & Sanskrit \emph{dvitīya} \\
\textbf{\pa{ਤੀਜਾ}} & \textbf{-\pa{ਜਾ}} & Sanskrit \emph{tṛtīya} \\
\textbf{\pa{ਚੌਥਾ}} & \textbf{-\pa{ਥਾ}} & Sanskrit \emph{caturtha} \\
\textbf{\pa{ਪੰਜਵਾਂ}} & \textbf{-\pa{ਵਾਂ}} & Sanskrit \emph{pañcama} \\
\bottomrule
\end{tabularx}

Every row of that last column is a whole word handed down, not a sum. So the useful question about this set is not \emph{which ones were built} — none of them were. It is \textbf{which ending got a second life}, and the answer is the one at the bottom: \textbf{-\pa{ਵਾਂ}} is the only one a Punjabi speaker can put on a number they have never made an ordinal of before.
\end{grammarlens}

\begin{tcolorbox}[breakable,colback=teal!4,colframe=teal!35!black,title={Before you move on}]
How many of the five were built out of their numbers? (\textbf{None} — \textbf{all five were inherited}.) Which ending can still be used on a new number? (\textbf{-\pa{ਵਾਂ}}.)
\end{tcolorbox}
